% TOP_PROBLEM: {"id": 3072, "title": "Edge-antipodal colorings of cubes", "category_id": 2, "category": {"id": 2, "name": "combinatorics", "display_name": "Combinatorics", "description": "Counting problems, graph theory, discrete structures.", "slug": "combinatorics", "order_index": 2, "created_at": "2026-07-31T15:26:25.671Z"}, "status": "open", "problem_number": "OPG-2359", "set_id": 12}
% FIRST_POSTED: 2026-10-01
% ATTEMPT_STATUS: unresolved
% UnsolvedMath ID 3072; source catalogue code OPG-2359
% !TeX program = lualatex
% GPT-6 Astra Ultra research attempt; independent partial work, 2026-10-01.
\documentclass[11pt,a4paper]{article}
\usepackage[margin=25mm]{geometry}
\usepackage{fontspec}
\setmainfont{lmroman10-regular.otf}[BoldFont=lmroman10-bold.otf,
 ItalicFont=lmroman10-italic.otf,BoldItalicFont=lmroman10-bolditalic.otf]
\usepackage{amsmath,amssymb,mathtools}
\usepackage{unicode-math}
\setmathfont{latinmodern-math.otf}
\AtBeginDocument{\let\setminus\smallsetminus}
\usepackage{xurl,hyperref}
\hypersetup{colorlinks=true,urlcolor=blue}
\setcounter{secnumdepth}{0}
\setlength{\parindent}{0pt}
\setlength{\parskip}{5pt}
\setlength{\emergencystretch}{3em}
\begin{document}
\section*{Edge-antipodal colorings of cubes}\label{3072}
{\small UnsolvedMath ID 3072 · OPG-2359 · Combinatorics · OpenGarden}
\subsection{Definitions and mathematical statement}
For $d\ge2$, let $Q_d$ be the graph on $\{0,1\}^d$ whose edges join
strings differing in exactly one coordinate. For a vertex
$x=(x_1,\ldots,x_d)$, define its antipode by
$\bar x=(1-x_1,\ldots,1-x_d)$. The antipodal edge to an unordered edge
$e=\{x,y\}$ is $\bar e=\{\bar x,\bar y\}$, which is again an edge.

An edge-antipodal two-colouring is a map
$\varphi:E(Q_d)\to\{0,1\}$ satisfying
$\varphi(\bar e)=1-\varphi(e)$ for every edge $e$.
No proper-colouring condition is imposed at a vertex.
The conjecture asserts that for every dimension and every such colouring
there are a vertex $x$ and a path from $x$ to $\bar x$ all of whose edges
have the same colour.

More explicitly, one seeks distinct vertices
$v_0,\ldots,v_t$ with $v_0=x$, $v_t=\bar x$, adjacent consecutive
vertices, and one $b\in\{0,1\}$ such that
$\varphi(\{v_{i-1},v_i\})=b$ for all $1\le i\le t$.
There is no requirement that $t=d$, that the path be shortest, or that
it change each coordinate just once.

Equivalently, some connected component of one of the two monochromatic
spanning subgraphs must contain an antipodal vertex pair. The pair is
chosen after the colouring and need not be a prescribed pair. Antipodality
is taken in the full $d$-cube; restricting the colouring to a lower-dimensional
face need not preserve this hypothesis.
\subsection{Short English statement}
Must a cube edge-colouring that reverses colour on antipodal edges contain a monochromatic path between some antipodal vertices?
\subsection{Research attempt: one-switch paths, components, and a status qualification}
\paragraph{Literature check and scope.}
GPT-6 Astra Ultra, 1 October 2026. The statement above retains the
catalogue's historical wording. Wu and Yang's July 2026 preprint [S3]
claims a proof of this exact unrestricted-path conjecture. Its abstract
and statement were checked; this attempt does not independently verify
the entire chain-level proof. Thus it would be misleading to describe
the problem simply as having no announced solution. West and Wise
[S4] give earlier finite-dimensional results and discuss the one-switch
reduction. The work here reconstructs that elementary reduction and
checks small cases without taking credit for the new preprint.

\paragraph{1. A path with one colour change suffices.}
Suppose a path from $x$ to $\bar x$ is red from $x$ to a vertex $y$
and blue from $y$ to $\bar x$. Apply the antipodal map to its blue
part. All those edges become red, giving a red path from $\bar y$
to $x$. Concatenate it with the original red part from $x$ to $y$.
This is a red walk from $\bar y$ to $y$. Repeated vertices can be
removed by deleting the intervening closed subwalk, so it contains
a red simple path between these antipodes. If either part was empty,
the original path was already monochromatic. Interchanging the colour
names handles the other order of the two parts.

Therefore any proof that every colouring under consideration has an
antipodal path with at most one colour change would prove the target.
No geodesic assumption was needed in this reduction. Conversely a
monochromatic antipodal path is already a path with zero changes, so
for antipodal colourings the two existence assertions are equivalent.

\paragraph{2. Component constraints are useful but not sufficient.}
Let $C$ be a red connected component. Its antipodal image $\bar C$
is a blue connected component: reflection preserves graph adjacency
and reverses colour, and maximality transfers by reflecting any larger
connected set back. If $|C|>2^{d-1}$, then $C$ contains both endpoints
of some antipodal pair, by placing at most one vertex in each of the
$2^{d-1}$ pairs. In a counterexample every monochromatic component
must therefore have size at most $2^{d-1}$, and every such component
is disjoint from its own antipodal image. These restrictions do not
imply existence of a large component. Counting total red and blue
edges only gives equal edge totals, not a lower bound above half the
vertices for one component.

\paragraph{3. Dimension two and a failed induction.}
The square has two antipodal pairs of edges. Each pair contributes
one red and one blue edge. The two red edges cannot be opposite,
because opposite edges must have different colours. They therefore
share a vertex and form a red path of length two whose endpoints
are antipodal. This proves the $d=2$ case directly.

For induction one might restrict to the face $x_1=0$ and invoke the
same assertion in dimension $d-1$. This is invalid: the antipodal
mate in $Q_d$ of an edge in that face lies in $x_1=1$. The original
condition imposes no relation between edges paired by the internal
antipodal map of the first face. In particular an arbitrary colouring
of its internal edges can be completed on the opposite face by
reflection with reversed colours. Thus the induction hypothesis is
not available on the chosen face.

\paragraph{Verification and exact remaining task.}
A direct exhaustive check represented vertices as bit strings, chose
one colour for each antipodal edge pair, and searched red components
for a pair $v,v\mathbin{\mathrm{xor}}(2^d-1)$. There are
$d2^{d-2}$ independent edge-pair choices. All $4$, $64$, and $65536$
colourings for $d=2,3,4$ passed. Testing red alone is sufficient:
reflection turns any blue antipodal path into a red one. These finite
checks confirm the conventions; they cannot establish arbitrary $d$.
The independent work here still lacks a general existence argument
for a one-switch path, or a complete verification of [S3]'s alternative
chain obstruction.

\paragraph{Reproduction of finite checks.}
The finite checks stated above were run with Python 3 using the repository
script [C1]. They verify the stated examples and finite ranges only.

\subsection{Final status}
UNRESOLVED as an independent full proof attempt. The elementary
reduction and small cases are verified; a claimed full solution is
explicitly credited to [S3]. This file does not reclassify that claim
as an original model solution or certify its entire proof.

\subsection{Sources}
\begin{itemize}
\item[S1] ULAM.AI and the UnsolvedMath Contributors, supplied problems.json, record 3072 (OPG-2359), OpenGarden collection. Catalogue identity and source transcription; CC BY 4.0.
\url{https://huggingface.co/datasets/ulamai/UnsolvedMath}.
\item[S2] Open Problem Garden, Edge-antipodal colorings of cubes. Original problem and discussion; the mathematical formulation below is expanded from the supplied source record.
\url{http://www.openproblemgarden.org/op/edge_antipodal_colorings_of_cubes}.
\item[S3] H. Wu and N. Yang, \emph{A Chain-Level Borsuk--Ulam Obstruction Proof of Norine's Antipodal-Coloring Conjecture}, arXiv:2607.19276v1, 21 July 2026. Source of the announced-proof qualification; checked 1 October 2026.
\url{https://arxiv.org/html/2607.19276v1}.
\item[S4] D. B. West and J. I. Wise, \emph{Antipodal edge-colorings of hypercubes}, author manuscript revised July 2017. Checked 1 October 2026.
\url{https://dwest.web.illinois.edu/pubs/antip.pdf}.
\item[C1] Reproducible finite-check code, executed 1 October 2026 with Python 3 (standard library only):
\path{scripts/check_attempt_batch_geometry_2026_10_01.py}.
\end{itemize}
\end{document}
